\section{Model architectures and training settings}
\label{app:spaces}
The model families represent different assumptions about how waveform
information can correct LED timing. Their detector-pair formulations are
described below; Tables~\ref{tab:space-neural-summary} and
\ref{tab:space-ridge-minirocket-summary} collect the configured search spaces.

\subsection{Linear ridge regression}
Linear models provide a computationally inexpensive baseline for testing
whether a weighted combination of waveform samples is sufficient for timing
correction. The direct estimator uses independent weights on the concatenated
detector waveforms and an intercept. The shared estimator instead applies
common weights to the sample-wise waveform difference, without an intercept,
enforcing detector-exchange antisymmetry:
\begin{equation}
h(s_1,s_2)=\mathbf{w}^{T}(s_1-s_2)
=\mathbf{w}^{T}s_1-\mathbf{w}^{T}s_2.
\end{equation}
Ridge regularization limits sensitivity to highly correlated samples. Its
strength is selected by internal leave-one-out cross-validation minimizing
mean squared error (MSE), rather than the outer CTR-based selection.

\subsection{Convolutional neural networks}
Convolutional neural networks (CNNs) are established waveform-timing models,
including for LED correction~\cite{onishi2022,loignon2025}. Their temporal
filters exploit local waveform structure with weights shared across sample
positions. The Onishi-inspired direct CNN combines the detector dimension in
the first convolution and predicts a joint correction. The independent CNN
uses separately parameterized detector scorers, whereas the shared CNN uses
one common scorer for both detector signals.

\subsection{Multilayer perceptrons}
The direct dense MLP performs nonlinear regression on the concatenated
waveforms. The shared MLP applies the same nonlinear scorer to both detector
signals and subtracts the two outputs,
\begin{equation}
h_\theta(s_1,s_2)=g_\theta(s_1)-g_\theta(s_2).
\end{equation}
The shared locally connected MLP preserves temporal neighborhoods without
sharing filters across positions. This is useful for LED-aligned waveforms,
where temporal position itself is informative, and provides an intermediate
assumption between dense and convolutional processing~\cite{elsayed2020}.

\subsection{MiniRocket feature regression}
MiniRocket is an efficient, largely deterministic convolutional feature
transformation developed for time-series classification~\cite{minirocket2021}.
Here its features are used for timing-error regression with ridge
regularization. The direct model transforms the waveform pair jointly; the
shared model applies the same univariate transform to both detector signals
and regresses on their feature difference.

\subsection{Training and configured spaces}
Dense MLPs, the locally connected MLP and temporal CNNs use at most 300
epochs, an internal holdout fraction of 0.2, loss-stop patience of 10 epochs,
a minimum loss improvement of $0.1\,\mathrm{ps}$, gradient clipping at 10,
and a gradient-stop threshold of $10^{-6}$ with patience 4. The Onishi-inspired
CNN instead uses a fixed 100-epoch schedule with learning-rate reductions at
epochs 30 and 60. Hyperparameters for neural and MiniRocket models are selected
from the training acquisition using mean held-out CTR in the three common
outer folds; the two linear ridge estimators use internal leave-one-out
cross-validation on MSE.

\subsection{Software libraries}
Table~\ref{tab:model-libraries} summarizes the principal scientific software
and declared minimum versions.
\begin{table}[tbp]
\centering
\scriptsize
\begin{tabularx}{\columnwidth}{@{}lclX@{}}
\toprule
Library & Min. & Role \\
\midrule
PyTorch & 2.2 & Neural models \\
scikit-learn & 1.4 & Ridge CV and scaling \\
sktime & 0.35 & MiniRocket transforms \\
NumPy & 1.26 & Waveform and bootstrap operations \\
SciPy & 1.11 & Numerical routines and Gaussian fits \\
Matplotlib & 3.8 & Figures \\
pandas & 2.1 & Tabular data and transformed features \\
Optuna & 4.0 & Optional adaptive parameter proposal \\
\bottomrule
\end{tabularx}
\caption{Principal scientific software and declared minimum versions.}
\label{tab:model-libraries}
\end{table}



\begin{table*}[tbp]
\centering
\footnotesize
\begin{tabularx}{\textwidth}{@{}p{0.08\textwidth}Xp{0.21\textwidth}p{0.12\textwidth}p{0.18\textwidth}@{}}
\toprule
Code & Architecture & Learning rate & Batch size & Optimizer / schedule \\
\midrule
S-MLP
& $\{[16],[128],[128,64],[16,16,16]\}$
& $\{10^{-2},10^{-3},10^{-4},10^{-5}\}$
& 32
& Nesterov SGD \\

D-MLP
& $\{[16],[128],[128,64],[16,16,16]\}$
& $\{10^{-2},10^{-3},10^{-4},10^{-5}\}$
& 32
& Adam \\

S-LCM
& Dense [128]; local kernels $(3,3)$, strides $(1,1)$; correction bound $500\,\mathrm{ps}$
& $\{10^{-2},10^{-3},10^{-4},10^{-5}\}$
& $\{32,128\}$
& Adam \\

S-CNN
& Channels [16,32]; kernels [5,3]; dense 64
& $\{10^{-3},10^{-4},10^{-5}\}$
& $\{32,128\}$
& Adam \\

D-CNN
& Channels [16,32]; kernels [5,3]; dense 64
& $\{10^{-3},10^{-4},10^{-5}\}$
& $\{32,128\}$
& Adam \\

D-ONI
& Channels [32,64]; kernels [5,3]; dense 256
& $\{10^{-3},10^{-4},10^{-5}\}$
& $\{32,128\}$
& Adam; 100 epochs, LR drops at 30/60 \\
\bottomrule
\end{tabularx}
\caption{Configured neural-model search spaces. Model codes are defined in
Table~\ref{tab:model-inventory}. All models use ReLU activations; training
settings shared across models are given in the text.}
\label{tab:space-neural-summary}
\end{table*}
\begin{table*}[tbp]
\centering
\scriptsize
\begin{tabularx}{\textwidth}{@{}llXXl@{}}
\toprule
Code & Estimator & Feature / regularization space & Selection & Intercept \\
\midrule
S-LIN & Shared linear ridge & 50 log-spaced ridge strengths from $10^{-8}$ to $10^{3}$ & Internal leave-one-out CV; MSE & None \\
D-LIN & Direct linear ridge & 50 log-spaced ridge strengths from $10^{-8}$ to $10^{3}$ & Internal leave-one-out CV; MSE & Included \\
S-MR & Shared MiniRocket + ridge & 10000 requested features; ridge $\alpha\in\{0.1,0.3,1,3,10,30,100,300,10^3,3\times10^3,10^4,3\times10^4,10^5,3\times10^5,10^6\}$ & Outer CV; mean CTR & None \\
D-MR & Direct MiniRocket + ridge & 10000 requested features; same ridge grid & Outer CV; mean CTR & Included \\
\bottomrule
\end{tabularx}
\caption{Configured linear and MiniRocket spaces. MiniRocket feature transforms
and scaling are fixed before supervised fitting.}
\label{tab:space-ridge-minirocket-summary}
\end{table*}


